\documentclass[11pt,a4paper]{article}
\usepackage{xeCJK}
\setCJKmainfont{Yu Mincho}
\setCJKsansfont{Yu Gothic}
\setCJKmonofont{Yu Gothic}
\usepackage[margin=25mm]{geometry}
\usepackage{amsmath,amssymb,amsthm,bm}
\usepackage{hyperref}

\newtheorem{assumption}{仮定}
\newtheorem{proposition}{命題}
\newtheorem{theorem}{定理}
\newtheorem{corollary}{系}
\newtheorem{remark}{注意}

\newcommand{\E}{\mathbb E}
\newcommand{\Prb}{\mathbb P}
\newcommand{\1}{\mathbf 1}
\newcommand{\Y}{\bm Y}
\newcommand{\D}{\bm D}
\newcommand{\indep}{\mathrel{\perp\!\!\!\perp}}

\title{欠測IV識別のベクトル化：\\
d'Haultfoeuille (2010), Proposition 2.1 and Theorem 2.3 の再構築}
\author{}
\date{}

\begin{document}
\maketitle

\begin{abstract}
d'Haultfoeuille (2010) のスカラー選択モデルを，項目別欠測指標
$\D$ と欠測変数ベクトル $\Y$ を持つ設定へ拡張する。中心となるのは
完全ケース指標 $R$ である。条件付き独立性と vector
$\mathcal B_{\Y}$-completeness の下で，完全ケース確率
$\pi(\Y)$，同時分布 $P(\Y,Z)$，および $P(R,\Y,Z)$ が識別される。
各非完全パターン $\D=\bm d\ne\bm 1_m$ の分布は，中心仮定だけからは
識別されないことも明確にする。
\end{abstract}

\section{原論文との記法対応}

原論文と本ノートの役割は次のように対応する。
\begin{center}
\begin{tabular}{lll}
\hline
対象 & d'Haultfoeuille (2010) & 本ノート \\
\hline
選択・完全ケース & $D\in\{0,1\}$ & $R=\1\{\D=\bm 1_m\}$ \\
欠測変数 & $Y$ & $\Y=(Y_1,\ldots,Y_m)'$ \\
常時観測IV & $Z$ & $Z$ \\
選択確率 & $P(y)$ & $\pi(\bm y)$ \\
\hline
\end{tabular}
\end{center}

項目別観測指標を $\D=(D_1,\ldots,D_m)'$ とし，
\begin{equation}
 R=\1\{\D=\bm 1_m\}
\end{equation}
を完全ケース指標とする。$Z$ は常に観測される。中心証明が利用する
観測分布は $P(Z)$ と完全ケースsubdistribution
\begin{equation}
 \Prb\{R=1,\Y\in d\bm y,Z\in dz\}
\end{equation}
である。$R=0$ における部分観測値は，識別を強める可能性はあるが，
以下の定理には用いない。

\section{構造から導かれる条件付き独立性}

\begin{assumption}[ベクトル構造]
ある可測関数 $\phi,\psi$ が存在して
\begin{equation}
 \Y=\phi(Z,\bm\varepsilon),
 \qquad
 \D=\psi(\Y,\bm\nu),
 \qquad
 \bm\nu\indep(Z,\bm\varepsilon)
 \label{eq:vector-structure}
\end{equation}
が成り立つ。
\end{assumption}

\begin{proposition}[ベクトル版 Proposition 2.1]
\label{prop:vector-ci}
仮定1の下で
\begin{equation}
 \D\indep Z\mid\Y,
 \qquad
 R\indep Z\mid\Y
 \label{eq:vector-ci}
\end{equation}
が成り立つ。
\end{proposition}

\begin{proof}
任意の欠測パターン $\bm d$ と値 $\bm y$ に対して
\begin{equation}
 A_{\bm y,\bm d}
 =\{\bm u:\psi(\bm y,\bm u)=\bm d\}
\end{equation}
と置く。$\Y$ は $(Z,\bm\varepsilon)$ の可測関数であり，
$\bm\nu\indep(Z,\bm\varepsilon)$ だから
$\bm\nu\indep(\Y,Z)$ である。したがって
\begin{align}
 \Prb(\D=\bm d\mid\Y=\bm y,Z=z)
 &=\Prb(\bm\nu\in A_{\bm y,\bm d}mid\Y=\bm y,Z=z)\\
 &=\Prb(\bm\nu\in A_{\bm y,\bm d})\\
 &=\Prb(\D=\bm d\mid\Y=\bm y).
\end{align}
よって $\D\indep Z\mid\Y$ である。$R$ は $\D$ の可測関数なので，
$R\indep Z\mid\Y$ も従う。
\end{proof}

\section{識別仮定}

完全ケース確率を
\begin{equation}
 \pi(\bm y)=\Prb(R=1\mid\Y=\bm y)
\end{equation}
と定義する。また，原論文と同じ関数クラス
\begin{equation}
 \mathcal B_{\Y}
 =\left\{a:\E|a(\Y)|<\infty,
 \ a(\Y)\text{ は下に有界 a.s.}\right\}
 \label{eq:admissible-class}
\end{equation}
を用いる。

\begin{assumption}[vector $\mathcal B_{\Y}$-completeness]
任意の $a\in\mathcal B_{\Y}$ について
\begin{equation}
 \E\{a(\Y)\mid Z\}=0\quad\text{a.s.}
 \quad\Longrightarrow\quad
 a(\Y)=0\quad\text{a.s.}
 \label{eq:vector-completeness}
\end{equation}
が成り立つ。
\end{assumption}

\begin{assumption}[positivity]
\begin{equation}
 \pi(\Y)>0\quad\text{a.s.}
 \label{eq:positivity}
\end{equation}
が成り立つ。
\end{assumption}

一様な下限 $\pi(\Y)\ge c>0$ は推定の安定性には有用であるが，以下の
識別証明には必要ない。

\section{完全ケース確率の一意識別}

\begin{theorem}[ベクトル版 Theorem 2.3]
\label{thm:vector-identification}
命題\ref{prop:vector-ci}の条件，vector
$\mathcal B_{\Y}$-completeness，および positivity が成り立つとする。
このとき $\pi(\Y)$，$P(\Y,Z)$，および $P(R,\Y,Z)$ は観測分布から
識別される。さらに $\E\|\Y\|<\infty$ なら $\E(\Y)$ も識別される。
\end{theorem}

\begin{proof}
証明を六段階に分ける。

\paragraph{1. 真の完全ケース確率がモーメント方程式を満たすこと.}
命題\ref{prop:vector-ci}より
\begin{align}
 \E\left\{\frac{R}{\pi(\Y)}\middle|Z\right\}
 &=\E\left[\left.
   \frac{\E(R\mid\Y,Z)}{\pi(\Y)}\right|Z\right]\\
 &=\E\left\{\left.
   \frac{\pi(\Y)}{\pi(\Y)}\right|Z\right\}=1.
 \label{eq:true-moment}
\end{align}

\paragraph{2. 候補モーメントの観測可能性.}
任意の可測な候補 $q:\mathbb R^m\to(0,1]$ に対し
\begin{equation}
 \E\left\{\frac{R}{q(\Y)}\middle|Z\right\}=1
 \label{eq:candidate-moment}
\end{equation}
を考える。$R/q(\Y)$ は $R=0$ ならゼロであり，$R=1$ のときだけ
$\Y$ を評価する。したがって式\eqref{eq:candidate-moment}の左辺は
観測分布から決まる。

\paragraph{3. 候補解から完備性を適用できる関数を作ること.}
$q$ が式\eqref{eq:candidate-moment}の解であるとし
\begin{equation}
 a(\Y)=\frac{\pi(\Y)}{q(\Y)}-1
 \label{eq:ratio-function}
\end{equation}
と置く。$\pi>0$ と $q>0$ より $a(\Y)\ge-1$ である。また
式\eqref{eq:candidate-moment}を $Z$ について積分すると
\begin{equation}
 \E\left\{\frac{\pi(\Y)}{q(\Y)}\right\}
 =\E\left\{\frac{R}{q(\Y)}\right\}=1.
\end{equation}
したがって
\begin{equation}
 \E|a(\Y)|
 \le \E\left\{\frac{\pi(\Y)}{q(\Y)}\right\}+1=2,
 \label{eq:ratio-integrability}
\end{equation}
ゆえに $a\in\mathcal B_{\Y}$ である。

\paragraph{4. vector completenessによる一意性.}
条件付き独立性と候補モーメントから
\begin{align}
 0
 &=\E\left\{\frac{R}{q(\Y)}-1\middle|Z\right\}\\
 &=\E\left\{\frac{\pi(\Y)}{q(\Y)}-1\middle|Z\right\}
 =\E\{a(\Y)\mid Z\}.
\end{align}
式\eqref{eq:vector-completeness}より $a(\Y)=0$ a.s.，したがって
$q(\Y)=\pi(\Y)$ a.s. である。よってモーメント方程式の解は一意で，
$\pi$ は識別される。

\paragraph{5. 母集団同時分布の復元.}
任意の可測集合 $A\subseteq\mathbb R^m\times\mathcal Z$ に対して
\begin{align}
 \Prb\{(\Y,Z)\in A\}
 &=\E\left[
 \frac{R}{\pi(\Y)}\1\{(\Y,Z)\in A\}
 \right].
 \label{eq:joint-recovery}
\end{align}
右辺は完全ケースsubdistributionと識別された $\pi$ だけから決まる。
同値な測度表示は
\begin{equation}
 P_{\Y,Z}(d\bm y,dz)
 =\frac{P(R=1,\Y\in d\bm y,Z\in dz)}{\pi(\bm y)}.
 \label{eq:measure-recovery}
\end{equation}
したがって $P(\Y,Z)$ が識別される。

\paragraph{6. $R=0$ 側と完全な三変量分布の復元.}
同じ集合 $A$ について
\begin{align}
 &\Prb\{R=0,(\Y,Z)\in A\}\
 &\qquad=
 \E\left[
 R\frac{1-\pi(\Y)}{\pi(\Y)}
 \1\{(\Y,Z)\in A\}
 \right].
 \label{eq:noncomplete-recovery}
\end{align}
$R=1$ 側は観測されているので，式\eqref{eq:noncomplete-recovery}と
合わせて $P(R,\Y,Z)$ 全体が識別される。
\end{proof}

\begin{corollary}[関数と平均の復元]
\label{cor:functional-recovery}
期待値が存在する任意の可測関数 $h$ について
\begin{equation}
 \E\{h(\Y,Z)\}
 =\E\left\{\frac{R}{\pi(\Y)}h(\Y,Z)\right\}.
 \end{equation}
特に $\E\|\Y\|<\infty$ なら
\begin{equation}
 \E(\Y)=\E\left\{\frac{R\Y}{\pi(\Y)}\right\}.
\end{equation}
\end{corollary}

\begin{remark}[識別されない対象]
中心仮定は $R\indep Z\mid\Y$ であるため，直接識別される欠測状態は
二値指標 $R$ である。各 $\bm d\ne\bm 1_m$ について
$P(\D=\bm d,\Y,Z)$ を識別するには，パターン別確率
$\Prb(\D=\bm d\mid\Y)$ に対する追加の除外制約，positivity，および
対応するモーメント方程式が必要である。したがって，本定理は
$P(R,\Y,Z)$ の識別を $P(\D,\Y,Z)$ の識別と読み替えない。
\end{remark}

\begin{remark}[項目別分解を仮定しないこと]
$Y_1,\ldots,Y_m$ は相関してよく，$\pi(\bm y)$ は複数成分へ同時に
依存してよい。証明は $\pi(\bm y)=\prod_j\pi_j(y_j)$ も
$D_j\indep D_k\mid\Y$ も用いない。
\end{remark}

\section{主論文の記法への写像}

本ノートはonepage版との連続性のため常時観測変数を $Z$ と書いた。
主論文では次の置換だけを行う。
\begin{equation}
 Z\longmapsto V,
 \qquad
 \E\{a(\Y)\mid Z\}=0
 \longmapsto
 \E\{a(\Y)\mid V\}=0.
\end{equation}
証明の各段階は変わらない。潜在構造
$F=g(V)+U$ と $\Y=\bm h(F)+\bm\varepsilon$ の識別は，ここで復元した
$P(\Y,V)$ を入力とする第二段階であり，本定理そのものとは区別する。

\begin{thebibliography}{9}
\bibitem{dhaultfoeuille2010}
d'Haultfoeuille, X. (2010).
A new instrumental method for dealing with endogenous selection.
\textit{Journal of Econometrics}, \textit{154}(1), 1--15.
\end{thebibliography}

\end{document}
